% =====================================================================
%  Práctica 2 - Sistema extendido de monitorización y control de temperatura
%  Sistemas Empotrados y Ubicuos - Grupo 4
%
%  Compilar: pdflatex -> pdflatex   (o: latexmk -pdf memoria.tex)
% =====================================================================
\documentclass{upo-master}

% ---------- Datos del documento ----------
\asignatura{Sistemas Empotrados y Ubicuos}
\tipodocumento{Práctica 2}
\title{Sistema extendido de monitorización y control de temperatura}
\subtitulo{Simulación en Proteus con Arduino (ATmega328P)}
\author{\texorpdfstring{\begin{tabular}[t]{@{}l@{}}Javier Gómez\\Daniel Jurado\\Manuel Ramos\end{tabular}}%
                       {Javier Gómez, Daniel Jurado, Manuel Ramos}}
\grupo{Grupo 4}
\tituloCorto{P2 - Control de temperatura}

% Estilo para el código Arduino (C++)
\lstdefinestyle{arduino}{language=C++, numbers=left,
  morekeywords={uint8_t, byte, boolean, HIGH, LOW, OUTPUT, INPUT},
  % Mismos acentos que la clase + el símbolo de grado (º) de los comentarios
  literate={á}{{\'a}}1 {é}{{\'e}}1 {í}{{\'i}}1 {ó}{{\'o}}1 {ú}{{\'u}}1
           {Á}{{\'A}}1 {É}{{\'E}}1 {Í}{{\'I}}1 {Ó}{{\'O}}1 {Ú}{{\'U}}1
           {ñ}{{\~n}}1 {Ñ}{{\~N}}1 {ü}{{\"u}}1 {¿}{{?`}}1 {¡}{{!`}}1
           {º}{{\textdegree}}1 {°}{{\textdegree}}1}

% ---------- Uso de IA ----------
\herramientaIA{Claude (Anthropic)}{Claude Opus 5.5}{Apoyo en el diseño del conexionado del motor y del servo, en el código de las cuestiones C1 a C4 y en la redacción de la memoria.}
\usoIA{La IA se ha usado como apoyo para traducir el montaje real del enunciado
a un esquema de Proteus, revisar las capturas del circuito, escribir los
\emph{sketches} de C1 a C4 y redactar esta memoria. El montaje en Proteus y la
simulación los ha hecho el grupo, que ha revisado el código y el texto.}

% ---------- Declaración de uso de IA en plural (trabajo en grupo) ----------
% La clase la redacta para un único autor; se redefine aquí para los tres
% miembros del grupo, sin modificar la plantilla.
\makeatletter
\renewcommand{\declaracionIA}{%
  \ifupo@iadeclarada\else
  \global\upo@iadeclaradatrue
  \clearpage
  \phantomsection
  \section*{Declaración de uso de inteligencia artificial}
  \addcontentsline{toc}{section}{Declaración de uso de inteligencia artificial}
  En la elaboración de este trabajo se han utilizado las siguientes
  herramientas de inteligencia artificial generativa:

  \begin{center}
  \small
  \begin{tabularx}{\textwidth}{>{\raggedright\arraybackslash}p{0.2\textwidth}
                               >{\raggedright\arraybackslash}p{0.2\textwidth}
                               >{\raggedright\arraybackslash}X}
    \toprule
    \rowcolor{upoazul!12}\textbf{Herramienta} & \textbf{Modelo / versión} & \textbf{Finalidad} \\
    \midrule
    \upo@iafilas
    \bottomrule
  \end{tabularx}
  \end{center}

  \ifx\upo@iatexto\@empty\else\upo@iatexto\par\fi

  La IA se ha utilizado como herramienta de apoyo. Todo el contenido
  generado con su ayuda ha sido revisado, contrastado y validado por los
  autores, y los resultados experimentales, capturas, datos y código
  ejecutado corresponden al trabajo realizado por el propio grupo.

  \begin{tcolorbox}[colback=upoazul!5, colframe=upoazul, boxrule=0.6pt,
      arc=2pt, fonttitle=\bfseries, title={Declaración de responsabilidad}]
    Nosotros, \textbf{Javier Gómez}, \textbf{Daniel Jurado} y
    \textbf{Manuel Ramos}, integrantes del \upo@grupo{} de la asignatura
    \upo@asignatura, declaramos que este documento es fruto de nuestro
    trabajo, que hemos indicado de forma veraz y completa el uso que hemos
    hecho de herramientas de inteligencia artificial y que hemos citado
    debidamente todas las fuentes consultadas. Asumimos la \textbf{plena
    responsabilidad} sobre su contenido ---exactitud técnica, corrección del
    código y de los resultados, idoneidad de las fuentes y originalidad--- y
    sobre el cumplimiento de la normativa académica de la \upoUniversidad.

    \vspace{0.8em}
    \upo@fecha \hfill
    \begin{tabular}[t]{@{}l@{}}Firmado: Javier Gómez\\
      \phantom{Firmado: }Daniel Jurado\\
      \phantom{Firmado: }Manuel Ramos\end{tabular}
  \end{tcolorbox}
  \fi}
\makeatother

\begin{document}

\maketitle

% Pie de página: el grupo en lugar de los tres nombres apilados
\fancyfoot[L]{\small\color{upogris}Grupo 4 --- Gómez, Jurado, Ramos}

% ======================= INTRODUCCIÓN =======================
\section{Introducción}

En esta práctica se amplía el sistema de monitorización de temperatura de la
Práctica~1 para que, además de medir, pueda \emph{actuar} sobre la temperatura
de la sala. Para ello se incorporan dos actuadores: un motor de corriente
continua, que hace de ventilador, y un servomotor, que señaliza la acción de
control. Todo el montaje se realiza en el entorno de simulación Proteus a
partir del montaje real descrito en el enunciado.

% ======================= C1 =======================
\section{C1: Sistema actuador ``motor''}

\begin{enunciado}
Implementar en Proteus el sistema actuador motor y una secuencia de pruebas
que compruebe sus dos modos de funcionamiento:
\begin{itemize}
  \item \textbf{Activar (velocidad máxima) / Desactivar}: 2\,s apagado,
        2\,s al 100\,\%, 2\,s apagado.
  \item \textbf{Activar (punto de trabajo) / Desactivar}: 2\,s apagado,
        2\,s al 25\,\%, 2\,s al 50\,\%, 2\,s al 100\,\%, 2\,s apagado.
\end{itemize}
\end{enunciado}

\subsection{Montaje}

El Arduino no puede suministrar la corriente que necesita el motor, por lo que
se usa un transistor como interruptor y una batería de 9\,V independiente.
Los componentes añadidos respecto a la Práctica~1 son los de la
Tabla~\ref{tab:componentes}.

\begin{table}[htbp]
  \centering
  \caption{Componentes del sistema actuador motor en Proteus.}
  \label{tab:componentes}
  \begin{tabular}{lll}
    \toprule
    \textbf{Componente} & \textbf{Dispositivo Proteus} & \textbf{Configuración} \\
    \midrule
    Motor de corriente continua & \texttt{MOTOR}   & Tensión nominal 9\,V \\
    MOSFET de canal N           & \texttt{2N7000}  & --- \\
    Diodo de protección         & \texttt{DIODE}   & --- \\
    Batería                     & Generador DC     & 9\,V \\
    \bottomrule
  \end{tabular}
\end{table}

Las conexiones son las siguientes:
\begin{itemize}
  \item \textbf{Puerta (G)} del MOSFET al pin digital \textbf{5} del Arduino
        (\texttt{IO5}, PD5/OC0B).
  \item \textbf{Drenador (D)} al terminal negativo del motor y al ánodo del diodo.
  \item \textbf{Fuente (S)} a tierra, común con la del Arduino.
  \item \textbf{+9\,V} al terminal positivo del motor y al cátodo del diodo.
\end{itemize}

\begin{concepto}{MOSFET como interruptor}
El 2N7000 es un MOSFET de canal N colocado entre el motor y tierra
(\emph{low-side switch}). Cuando el pin del Arduino pone 5\,V en la puerta,
el transistor conduce entre D y S y circula corriente desde la batería a
través del motor. Con 0\,V en la puerta, el circuito queda abierto y el motor
se detiene. Como la puerta está aislada, el Arduino solo aporta tensión, no
corriente.
\end{concepto}

\begin{concepto}{Diodo de protección (\emph{flyback})}
El motor es una carga inductiva: al cortar su corriente de golpe genera un
pico de tensión que podría dañar el MOSFET. El diodo, colocado en paralelo con
el motor y polarizado en inverso, no conduce en funcionamiento normal; al
apagarse el transistor ofrece un camino para que la corriente de la bobina se
descargue, limitando la tensión del drenador a unos 0,7\,V por encima de la
alimentación.
\end{concepto}

\captura{c1-circuito.png}{Esquema del sistema actuador motor en Proteus.}

\subsection{Control de la velocidad: PWM}

El primer modo solo necesita encender o apagar el motor, así que basta con
\texttt{digitalWrite()}. Para el segundo modo hay que fijar una velocidad
intermedia. Como el ATmega328P no tiene salida analógica real,
\texttt{analogWrite()} genera una señal \textbf{PWM}: una onda cuadrada de
0--5\,V cuyo ciclo de trabajo es proporcional al valor escrito (0--255). El
motor responde a la tensión media, de modo que gira más rápido cuanto mayor es
el valor.

Solo los pines conectados a una salida de comparación de un temporizador
(\texttt{OCxx} en el símbolo del ATmega328P) pueden generar PWM: los
pines 3, 5, 6, 9, 10 y 11. Se eligió el \textbf{pin 5} (OC0B, Timer0) porque
está libre y porque el Timer1 (pines 9 y 10) lo usará la librería
\texttt{Servo} en la cuestión C2.

Para pasar del porcentaje de velocidad al valor PWM se redondea al entero más
cercano, de forma que el 25\,\%, 50\,\% y 100\,\% corresponden a 64, 128 y 255:

\begin{lstlisting}[style=arduino, caption={Conversión de porcentaje a valor PWM.}]
uint8_t porcentajeAPwm(uint8_t porcentaje)
{
  if (porcentaje > 100)
  {
    porcentaje = 100;
  }
  return (uint8_t)((porcentaje * (unsigned int)PWM_MAX + 50) / 100); // Se suma 50 para redondear
}
\end{lstlisting}

\subsection{Secuencia de pruebas}

El programa ejecuta en bucle las dos secuencias del enunciado, cada paso
durante 2\,s, y escribe cada cambio de estado por el puerto serie. El código
completo se incluye en el Anexo~\ref{anexo:c1}.

\begin{lstlisting}[style=arduino, caption={Secuencias de prueba de los dos modos.}]
void pruebaModoDigital()           // Modo 1: digitalWrite
{
  motorDigital(false); delay(DURACION_PASO_MS);
  motorDigital(true);  delay(DURACION_PASO_MS);
  motorDigital(false); delay(DURACION_PASO_MS);
}

void pruebaModoPwm()               // Modo 2: analogWrite (PWM)
{
  const uint8_t pasos[] = {0, 25, 50, 100, 0};
  for (uint8_t i = 0; i < sizeof(pasos); i++)
  {
    motorVelocidad(pasos[i]);
    delay(DURACION_PASO_MS);
  }
}

void loop()
{
  pruebaModoDigital();   // 6 s
  pruebaModoPwm();       // 10 s
}
\end{lstlisting}

\subsection{Resultado}

Al simular, el motor permanece parado, arranca a velocidad máxima y vuelve a
pararse en el primer modo; en el segundo, va aumentando su velocidad en los
escalones del 25\,\%, 50\,\% y 100\,\% antes de detenerse. La secuencia
completa dura 16\,s y se repite indefinidamente.

\begin{aviso}
Con \texttt{analogWrite()} en un pin sin PWM (por ejemplo, una entrada
analógica como A1) el Arduino no genera la señal: pone el pin a nivel bajo
para valores menores que 128 y a nivel alto para el resto, por lo que el
motor solo podría estar apagado o a máxima velocidad.
\end{aviso}

% ======================= C2 =======================
\section{C2: Sistema actuador ``servo''}

\begin{enunciado}
Implementar en Proteus el sistema actuador servo y una secuencia de pruebas
que lo lleve a los ángulos 0\textdegree, 45\textdegree, 90\textdegree,
45\textdegree{} y 0\textdegree, manteniendo cada posición durante 2\,s.
\end{enunciado}

\subsection{Montaje}

El servomotor se usará como indicador de la acción de control del ventilador:
0\textdegree{} con el motor apagado y 90\textdegree{} con el motor a máxima
velocidad. A diferencia del motor de C1, el servo tiene su propio circuito de
control interno, por lo que se conecta directamente al Arduino con tres cables.
Los componentes añadidos son los de la Tabla~\ref{tab:componentes-servo}.

\begin{table}[htbp]
  \centering
  \caption{Componentes del sistema actuador servo en Proteus.}
  \label{tab:componentes-servo}
  \begin{tabular}{lll}
    \toprule
    \textbf{Componente} & \textbf{Dispositivo Proteus} & \textbf{Configuración} \\
    \midrule
    Servomotor              & \texttt{MOTOR-PWMSERVO} & 0\textdegree--180\textdegree, pulsos de 544--2400\,\textmu s \\
    Condensador de desacoplo & \texttt{CAP-POL}       & 100\,\textmu F (\texttt{100uF}) \\
    \bottomrule
  \end{tabular}
\end{table}

Las propiedades del servo se configuran como indica el enunciado
(\emph{Minimum Angle} 0, \emph{Maximum Angle} +180, \emph{Rotational Speed} 50,
\emph{Minimum Control Pulse} 544u y \emph{Maximum Control Pulse} 2400u) para
que su comportamiento sea realista. Las conexiones son:
\begin{itemize}
  \item \textbf{Control} del servo al pin digital \textbf{9} del Arduino
        (\texttt{IO9}, PB1/OC1A).
  \item \textbf{Alimentación (+)} a los \textbf{5\,V} del Arduino, no a la
        batería de 9\,V del motor.
  \item \textbf{Tierra (--)} a GND.
  \item \textbf{Condensador de 100\,\textmu F} en paralelo con la alimentación
        del servo: terminal positivo a 5\,V y negativo a GND.
\end{itemize}

\begin{concepto}{Condensador de desacoplo}
Cuando el servo empieza a moverse pide un pico de corriente que podría hacer
caer la tensión de alimentación del Arduino. El condensador, colocado junto al
servo, actúa como una pequeña reserva de carga que suministra ese pico y
aísla al resto del circuito de él. Es electrolítico, por lo que tiene
polaridad: el terminal negativo debe ir siempre a tierra.
\end{concepto}

\captura{c2-circuito.png}{Esquema del sistema actuador servo en Proteus.}

\subsection{Control de la posición}

El servo no se controla por el nivel de tensión, sino por la \textbf{duración
de un pulso} que se repite cada 20\,ms: un pulso de 544\,\textmu s lo lleva a
0\textdegree{} y uno de 2400\,\textmu s a 180\textdegree. Su circuito interno
compara la posición del eje, medida con un potenciómetro, con la que indica el
pulso, y mueve el motor hasta que coinciden.

De generar estos pulsos se encarga la librería \texttt{Servo}: basta con
asociar el objeto al pin con \texttt{attach()} e indicar el ángulo con
\texttt{write()}. La librería usa el \textbf{Timer1}, lo que desactiva el PWM
de \texttt{analogWrite()} en los pines 9 y 10. Por eso el servo va en el pin 9
y el motor de C1 se dejó en el pin 5, que depende del Timer0, de forma que
ambos actuadores pueden funcionar a la vez en las siguientes cuestiones.

\subsection{Secuencia de pruebas}

El programa recorre en bucle la lista de ángulos del enunciado, manteniendo
cada uno 2\,s e informando por el puerto serie. El motor de C1 se deja apagado
durante la prueba. El código completo se incluye en el Anexo~\ref{anexo:c2}.

\begin{lstlisting}[style=arduino, caption={Configuración del servo y secuencia de prueba.}]
#include <Servo.h>

Servo servo;
const uint8_t PIN_SERVO = 9;

void setup()
{
  servo.attach(PIN_SERVO);   // Asocia el objeto al pin del servo
  servo.write(0);            // Posición inicial
}

void pruebaServo()
{
  const uint8_t angulos[] = {0, 45, 90, 45, 0};
  for (uint8_t i = 0; i < sizeof(angulos); i++)
  {
    moverServo(angulos[i]);  // servo.write() + mensaje por el serie
    delay(DURACION_PASO_MS);
  }
}
\end{lstlisting}

\subsection{Resultado}

Al simular, el servo parte de 0\textdegree, gira a 45\textdegree{} y a
90\textdegree, y vuelve a 45\textdegree{} y a 0\textdegree, permaneciendo 2\,s
en cada posición. El giro entre posiciones no es instantáneo, ya que depende
de la velocidad de rotación configurada en el componente. La secuencia dura
10\,s y se repite indefinidamente.

\subsection{Problema encontrado: rango del servo}

En la primera simulación el servo no alcanzaba los ángulos pedidos: al
ordenar 0\textdegree, 45\textdegree{} y 90\textdegree{} se colocaba en
$-90$\textdegree, $-88{,}5$\textdegree{} y $-5$\textdegree. El código era
correcto; el problema estaba en que el componente \texttt{MOTOR-PWMSERVO}
conservaba sus propiedades por defecto, con un rango distinto al que usa la
librería \texttt{Servo} (Tabla~\ref{tab:rango-servo}).

\begin{table}[htbp]
  \centering
  \caption{Rango de la librería \texttt{Servo} y rango por defecto del servo de Proteus.}
  \label{tab:rango-servo}
  \begin{tabular}{lcc}
    \toprule
    & \textbf{Librería \texttt{Servo}} & \textbf{Proteus (por defecto)} \\
    \midrule
    Ángulos       & 0\textdegree{} a 180\textdegree & $-90$\textdegree{} a $+90$\textdegree \\
    Pulsos        & 544 a 2400\,\textmu s           & 1000 a 2000\,\textmu s \\
    \bottomrule
  \end{tabular}
\end{table}

La librería traduce 0\textdegree, 45\textdegree{} y 90\textdegree{} en pulsos
de 544, 1008 y 1472\,\textmu s. El servo, con su rango por defecto,
interpretaba esos pulsos de otra forma:
\begin{itemize}
  \item \textbf{544\,\textmu s} queda por debajo de su mínimo de 1000\,\textmu s,
        así que se quedaba en el tope de $-90$\textdegree.
  \item \textbf{1008\,\textmu s} está apenas por encima del mínimo, lo que
        corresponde a $-88{,}5$\textdegree.
  \item \textbf{1472\,\textmu s} está cerca de la mitad del rango, lo que
        corresponde a unos $-5$\textdegree.
\end{itemize}

La solución fue configurar el componente con los parámetros del enunciado
(ángulo de 0 a +180, pulsos de \texttt{544u} a \texttt{2400u}), de modo que
librería y servo usen el mismo rango. Tras el cambio, el servo se coloca
correctamente en cada posición.

\begin{aviso}
En las propiedades de Proteus hay que escribir la unidad \texttt{u}
(\emph{micro}) en los pulsos: \texttt{544u} y \texttt{2400u}. Sin ella, los
valores se interpretan en segundos.
\end{aviso}

% ======================= C3 =======================
\section{C3: Sistema de monitorización y control de temperatura (base)}

\begin{enunciado}
Extender el sistema de monitorización de temperatura de la Práctica~1 con
los actuadores de C1 y C2:
\begin{itemize}
  \item Si la temperatura es \textbf{menor que el umbral} y no se pulsa el
        botón, el motor permanece desactivado y el servo en 0\textdegree{}
        (enfriamiento parado).
  \item Si la temperatura es \textbf{mayor o igual que el umbral}, o se pulsa
        el botón, el motor se activa al 100\,\% y el servo se coloca en
        90\textdegree{} (enfriamiento en marcha).
\end{itemize}
\end{enunciado}

\subsection{Montaje}

No se añade ningún componente nuevo: el circuito reúne en un mismo esquema
el termómetro de la Práctica~1 y los dos actuadores de C1 y C2
(Tabla~\ref{tab:pines-c3}).

\begin{table}[htbp]
  \centering
  \caption{Asignación de pines del sistema completo.}
  \label{tab:pines-c3}
  \begin{tabular}{lll}
    \toprule
    \textbf{Elemento} & \textbf{Pin} & \textbf{Función} \\
    \midrule
    Sensor LM35   & A0 & Entrada analógica (temperatura) \\
    Pulsador      & 4  & Entrada digital (forzar enfriamiento) \\
    LED rojo      & 3  & Enfriamiento en marcha \\
    LED verde     & 2  & Enfriamiento parado \\
    Motor (MOSFET)& 5  & Salida PWM (Timer0) \\
    Servo         & 9  & Pulsos de la librería \texttt{Servo} (Timer1) \\
    \bottomrule
  \end{tabular}
\end{table}

\captura{c3-circuito.png}{Esquema completo del sistema de monitorización y control en Proteus.}

\subsection{Lógica de control}

El programa parte del código de la Práctica~1, que ya leía el LM35 cada
segundo sin bloquear el bucle (comprobando \texttt{millis()} en lugar de usar
\texttt{delay()}) y atendía el pulsador en cada vuelta. Sobre esa base se
añade una única función, \texttt{aplicarEnfriamiento()}, que actúa a la vez
sobre los tres indicadores del estado del sistema:

\begin{lstlisting}[style=arduino, caption={Activación y desactivación del sistema de enfriamiento.}]
void aplicarEnfriamiento(bool activo)
{
  if (activo)
  {
    analogWrite(PIN_MOTOR, PWM_MOTOR_MAX);   // Motor al 100 %
    servo.write(ANGULO_MARCHA);              // Servo a 90º
    encenderLeds(true, false);               // LED rojo
  }
  else
  {
    analogWrite(PIN_MOTOR, 0);               // Motor parado
    servo.write(ANGULO_PARADO);              // Servo a 0º
    encenderLeds(false, true);               // LED verde
  }
}
\end{lstlisting}

En cada vuelta del bucle se decide si hay que enfriar (temperatura alta
\emph{o} botón pulsado) y solo se llama a esa función cuando la decisión
cambia, para no reenviar la misma orden a los actuadores continuamente:

\begin{lstlisting}[style=arduino, caption={Decisión de control en \texttt{loop()}.}]
bool botonPulsado = (digitalRead(PIN_BOTON) == HIGH);
bool debeEnfriar  = botonPulsado || temperaturaAlta;

if (debeEnfriar != enfriando)
{
  enfriando = debeEnfriar;
  aplicarEnfriamiento(enfriando);
}
\end{lstlisting}

El código completo se incluye en el Anexo~\ref{anexo:c3}.

\begin{nota}[Histéresis]
Como en la Práctica~1, la comparación con el umbral (24\,\textdegree C) usa
una histéresis de 0,5\,\textdegree C: el enfriamiento se activa al llegar a
24\,\textdegree C, pero no se desactiva hasta bajar de 23,5\,\textdegree C.
Con un control todo/nada esto es aún más importante que con los LEDs: si la
temperatura oscila justo en el umbral, sin histéresis el motor y el servo se
encenderían y apagarían continuamente. Si se quiere el comportamiento estricto
del enunciado basta con poner \texttt{HISTERESIS = 0}.
\end{nota}

\subsection{Resultado y comportamiento esperado}

En la simulación se modifica la temperatura del LM35 con sus flechas y se
comprueba que, al alcanzar el umbral, el motor arranca a máxima velocidad, el
servo gira a 90\textdegree{} y se enciende el LED rojo; al bajar de nuevo, todo
vuelve al estado de reposo. Mientras se mantiene pulsado el botón, el
enfriamiento se activa con independencia de la temperatura.

Como indica el enunciado, Proteus no dispone de un modelo térmico de la
sala, así que no puede verse el efecto del ventilador sobre la temperatura.
En un montaje real cabría esperar un comportamiento \textbf{todo/nada}: el
ventilador arrancaría a plena potencia al superar el umbral, la temperatura
bajaría hasta salir de la banda de histéresis, el ventilador se pararía y la
temperatura volvería a subir. El resultado sería una oscilación continua en
torno al umbral, con el motor alternando entre parado y máxima velocidad
aunque la diferencia de temperatura sea mínima. Esta limitación es la que
motiva el control proporcional de la cuestión C4.

% ======================= C4 =======================
\clearpage
\section{C4: Sistema de monitorización y control de temperatura (avanzado)}

\begin{enunciado}
Modificar C3 para que la acción de control (la velocidad del motor) sea
\textbf{proporcional a la diferencia} entre la temperatura real y la
temperatura umbral, probando las constantes de proporcionalidad 20, 50
y 100. El ángulo del servo debe señalizar el porcentaje de la velocidad
máxima, haciendo corresponder 0--100\,\% con 0--90\textdegree.
\end{enunciado}

\subsection{Ley de control}

El montaje es el mismo que en C3; solo cambia el programa. En lugar de
encender o apagar el motor, la velocidad se calcula a partir del
\textbf{error}, la diferencia entre la temperatura medida $T$ y el umbral
$T_u = 24$\,\textdegree C:
\begin{equation}
  u = K_p \,(T - T_u), \qquad
  \text{velocidad (\%)} = \min\bigl(\max(u,\,0),\,100\bigr)
\end{equation}

La acción se \textbf{satura} al rango físico del motor: por debajo del umbral
el error es negativo y el motor queda parado, y por encima de
$100/K_p$\,\textdegree C de diferencia el motor ya va al 100\,\%. Ese
porcentaje se traduce a los actuadores con dos conversiones lineales
redondeadas: al valor PWM (0--255), como en C1, y al ángulo del servo
(0--90\textdegree). El pulsador, como en C3, fuerza el 100\,\% mientras se
mantiene pulsado.

\begin{lstlisting}[style=arduino, caption={Cálculo de la acción de control proporcional.}]
const float UMBRAL = 24.0;
const float KP     = 20.0;   // Probar con 20, 50 y 100

uint8_t calcularAccionControl(float temperatura)
{
  float error  = temperatura - UMBRAL;   // > 0 si hace más calor del deseado
  float accion = KP * error;             // Acción proporcional al error

  if (accion <= 0.0)   return 0;         // Por debajo del umbral: no enfriar
  if (accion >= 100.0) return 100;       // Saturación: velocidad máxima
  return (uint8_t)(accion + 0.5);        // Redondeo al entero más cercano
}

void aplicarEnfriamiento(uint8_t porcentaje)
{
  analogWrite(PIN_MOTOR, porcentajeAPwm(porcentaje));      // 0..255
  servo.write(porcentajeAAngulo(porcentaje));              // 0..90º
  encenderLeds(porcentaje > 0, porcentaje == 0);
}
\end{lstlisting}

A diferencia de C3, aquí no se usa histéresis: el control ya es continuo, de
modo que cerca del umbral el motor gira despacio en lugar de conmutar entre
parado y máxima velocidad. El código completo se incluye en el
Anexo~\ref{anexo:c4}.

\subsection{Resultados con cada constante}

La Tabla~\ref{tab:kp} recoge la acción de control teórica para varias
temperaturas con cada valor de $K_p$, y el ángulo correspondiente del servo.
En la simulación los valores pueden diferir ligeramente por la resolución del
ADC (unos 0,49\,\textdegree C por paso).

\begin{table}[htbp]
  \centering
  \caption{Velocidad del motor (\%) y ángulo del servo según $K_p$ (umbral 24\,\textdegree C).}
  \label{tab:kp}
  \begin{tabular}{ccccc}
    \toprule
    \textbf{Temperatura} & \textbf{Error} & $\mathbf{K_p = 20}$ & $\mathbf{K_p = 50}$ & $\mathbf{K_p = 100}$ \\
    \midrule
    23\,\textdegree C & $-1$ & 0\,\% \;(0\textdegree)   & 0\,\% \;(0\textdegree)   & 0\,\% \;(0\textdegree) \\
    24\,\textdegree C & 0    & 0\,\% \;(0\textdegree)   & 0\,\% \;(0\textdegree)   & 0\,\% \;(0\textdegree) \\
    25\,\textdegree C & 1    & 20\,\% \;(18\textdegree) & 50\,\% \;(45\textdegree) & 100\,\% \;(90\textdegree) \\
    26\,\textdegree C & 2    & 40\,\% \;(36\textdegree) & 100\,\% \;(90\textdegree) & 100\,\% \;(90\textdegree) \\
    27\,\textdegree C & 3    & 60\,\% \;(54\textdegree) & 100\,\% \;(90\textdegree) & 100\,\% \;(90\textdegree) \\
    28\,\textdegree C & 4    & 80\,\% \;(72\textdegree) & 100\,\% \;(90\textdegree) & 100\,\% \;(90\textdegree) \\
    29\,\textdegree C & 5    & 100\,\% \;(90\textdegree) & 100\,\% \;(90\textdegree) & 100\,\% \;(90\textdegree) \\
    \midrule
    \multicolumn{2}{l}{\textbf{Saturación a partir de}} & $+5$\,\textdegree C & $+2$\,\textdegree C & $+1$\,\textdegree C \\
    \bottomrule
  \end{tabular}
\end{table}

\captura{c4-simulacion.png}{Simulación del control proporcional en Proteus.}

\subsection{Análisis de las constantes de proporcionalidad}

Igual que en C3, Proteus no permite ver el efecto real del ventilador sobre
la temperatura, pero sí razonar el comportamiento esperado:

\begin{itemize}
  \item \textbf{$K_p = 20$ (acción suave).} El motor recorre todo su rango a lo
        largo de 5\,\textdegree C, así que la velocidad se ajusta
        gradualmente y sin cambios bruscos, con menos ruido y consumo. Su
        inconveniente es una respuesta lenta y un \textbf{error permanente}
        mayor: un control proporcional solo actúa si hay error, de modo que
        la sala se estabilizaría por encima del umbral, en el punto en que la
        velocidad del ventilador compense el calor que entra. Además, a
        pocas décimas del umbral el motor recibe un PWM bajo con el que un
        motor real podría no llegar a arrancar (zona muerta).
  \item \textbf{$K_p = 50$ (compromiso).} Satura a 2\,\textdegree C de
        diferencia: reacciona con más fuerza y reduce el error permanente,
        manteniendo todavía un tramo proporcional útil.
  \item \textbf{$K_p = 100$ (acción agresiva).} Satura con solo
        1\,\textdegree C de error, por lo que el sistema se parece mucho al
        todo/nada de C3: reacción rápida y error permanente pequeño, pero con
        tendencia a oscilar y a forzar el motor. Además, cada paso del ADC
        (unos 0,49\,\textdegree C) cambia la acción en unos 49 puntos, así que
        el ruido de la medida se traduce en saltos grandes de velocidad.
\end{itemize}

En resumen, cuanto mayor es $K_p$, más rápida es la respuesta y menor el error
en régimen permanente, a costa de una actuación más brusca y más sensible al
ruido; cuanto menor es, más suave y estable es el control, pero más se aleja
la temperatura final del umbral. Para un ventilador doméstico, $K_p = 50$
parece el mejor equilibrio entre las tres.

% ======================= USO DE IA =======================
\declaracionIA

% ======================= ANEXO =======================
\appendix
\section{Código completo de C1}\label{anexo:c1}

\lstinputlisting[style=arduino, caption={\texttt{C1\_motor.ino}}]{../C1_motor.ino}

\section{Código completo de C2}\label{anexo:c2}

\lstinputlisting[style=arduino, caption={\texttt{C2\_servo.ino}}]{../C2_servo.ino}

\section{Código completo de C3}\label{anexo:c3}

\lstinputlisting[style=arduino, caption={\texttt{C3\_control.ino}}]{../C3_control.ino}

\section{Código completo de C4}\label{anexo:c4}

\lstinputlisting[style=arduino, caption={\texttt{C4\_proporcional.ino}}]{../C4_proporcional.ino}

\end{document}
